\section{Vector spaces}
\label{section-vector-spaces}

\noindent
One reason fields are so nice is that the theory of modules over fields
(i.e. vector spaces), is very simple.

\begin{lemma}
\label{lemma-vector-space-is-free}
If $k$ is a field, then every $k$-module is free.
\end{lemma}

\begin{proof}
Indeed, by linear algebra we know that a $k$-module (i.e. vector space)
$V$ has a {\it basis} $\mathcal{B} \subset V$, which defines an isomorphism
from the free vector space on $\mathcal{B}$ to $V$.
\end{proof}

\begin{lemma}
\label{lemma-field-semi-simple}
Every exact sequence of modules over a field splits.
\end{lemma}

\begin{proof}
This follows from Lemma \ref{lemma-vector-space-is-free} as every vector
space is a projective module.
\end{proof}

\noindent
This is another reason why much of the theory in future chapters will not say
very much about fields, since modules behave in such a simple manner.
Note that Lemma \ref{lemma-field-semi-simple} is a statement about the
{\it category} of $k$-modules (for $k$ a field), because the notion of
exactness is inherently arrow-theoretic, i.e., makes use of purely categorical
notions, and can in fact be phrased within a so-called {\it abelian category}.

\medskip\noindent
Henceforth, since the study of modules over a field is linear algebra, and
since the ideal theory of fields is not very interesting, we shall study what
this chapter is really about: {\it extensions} of fields.


